\chapter*{この報告書について}
\addcontentsline{toc}{chapter}{この報告書について}

\section*{目的}

LLM-BLM の研究では、言語モデルが決算発表文を読んだ点数を Black--Litterman モデルの見通しに入れて、ポートフォリオを作った。読み取りは発表をはさむ 2 日の株価の反応を強く説明したが、その情報は発表の当日にほぼ株価に入りきり、取引費用の後で時価総額加重を上回ることはできなかった。

では、プロのヘッジファンドは何をして収益を上げているのか。どんな戦略があり、どんなアルゴリズムを使い、その根拠になる研究は何を示していて、どこで失敗してきたのか。この報告書は、それを体系的にまとめたものである。読者には、金融の予備知識のない理系の大学生（線形代数・確率・統計・最適化の基本は知っている）を想定した。

\section*{構成}

\begin{description}[leftmargin=2em,style=nextline]
\item[第 I 部　基本（第 1〜11 章）]
  ヘッジファンドの仕組み、戦略の分類、成績の測り方、ファクター、シグナル、予測モデル、ポートフォリオ構築、リスク管理、取引コストと執行、バックテストと過学習、マーケットメイク。教科書や古典的な論文で確立していて、新しく調べなくても書ける標準的な内容である。式と小さな数値例で、仕組みが手で追えるように書いた。
\item[第 II 部　研究で分かっていること（第 12 章以降）]
  Gemini Deep Research に、切り口を変えた 3 つの依頼（戦略の種類から、作り方の工程から、成績の実証と失敗から）で文献を調べさせ、挙がった文献を照合したうえで、戦略ごと・工程ごとに整理した。第 I 部で説明した道具が、実際のデータでどれだけ効いたか、どこで崩れたかを扱う。最後の章で、この研究（LLM-BLM）への含みをまとめる。
\item[付録]
  用語の一覧と、文献の照合の結果。
\end{description}

\section*{文献の扱い}

\begin{itemize}
\item 第 I 部の引用は、分野の古典としてよく知られた論文と教科書に限った。全体の参考になる教科書として、Grinold・Kahn \cite{grinoldkahn2000}、Pedersen \cite{pedersen2015}、Ilmanen \cite{ilmanen2011}、López de Prado \cite{lopezdeprado2018}、Bouchaud ほか \cite{bouchaud2018} がある。
\item 第 II 部の文献は、Deep Research の 3 本の答え（\code{03\_方針と文献調査/ヘッジファンドの手法/}）に挙がったものと、それを補うために自分で選んだものである。報告書で引くすべての文献を、\code{refcheck.py} で題名から arXiv・OpenAlex・Crossref と照合し、実在が確かめられたものだけを使った。Deep Research の答えの誤りや、確かめられなかった記述は、付録~\ref{ch:verify} に一覧にした。
ef{ch:verify} に一覧にした。
\item 論文の数値は、その論文の標本（市場・期間・手数料や費用の前後）のものである。別の期間や市場で同じ値が出るとは限らない。本文では、できるだけ標本の範囲を一緒に書いた。
\item Deep Research の答えをそのまま写してはいない。数値は、照合で原典の要旨や表に当たれたもの、または分野でよく知られたものだけを書き、それ以外は「〜と報告されている」と出典を示して書くか、使わなかった。
\end{itemize}

\begin{notebox}
この報告書は勉強のための整理であり、投資の助言ではない。ヘッジファンドの成績のデータには、第 II 部で述べるように大きな偏りがある。論文で強かった戦略が、これからも同じように儲かるとは限らない。
\end{notebox}
